% Source: https://artofproblemsolving.com/community/c3h1518198_roots_integer
% Stable identity: aops-topic-1518198-q01
No. Suppose $1782^{12}+1841^{12}=N^{12}$ for an integer $N$, and work modulo the prime $13$.

By Fermat's little theorem, $a^{12}\equiv1\pmod{13}$ if $13\nmid a$, and $a^{12}\equiv0$ if $13\mid a$. So every twelfth power is $\equiv0$ or $1\pmod{13}$.

Now $1782=13\cdot137+1$ and $1841=13\cdot141+8$, so neither is divisible by $13$ and
\[
1782^{12}+1841^{12}\equiv1+1=2\pmod{13}.
\]
But $N^{12}\equiv0$ or $1$, never $2$. So no such $N$ exists, and the number is $\boxed{\text{not an integer}}$. $\blacksquare$

(By Fermat's Last Theorem no sum of two positive twelfth powers is a twelfth power; the argument above avoids that deep result. A calculator showing $1922$ is only an approximation: the true value is about $1921.99999996$.)
